\section{Introduction}
\label{sec:introduction}

Modern computational advertising is currently anchored in a passive, correlational paradigm. At the heart of the multibillion-dollar Real-Time Bidding (RTB) ecosystem lies the maximization of Click-Through Rate (CTR) and short-term Return on Ad Spend (ROAS), typically formulated as a point-estimation problem: $P(\text{click} | \text{ad}, \text{user})$. This approach, largely driven by Reinforcement Learning (RL) \citep{sutton2018reinforcement} and deep discriminative architectures \citep{he2016deep, devlin2018bert}, treats the user as a stochastic ``bandit arm'' to be sampled. While effective at capturing immediate user intent, this ``passive reward'' maximization ignores the fundamental agency of the human user and the long-term cognitive costs of engagement. The result is a system prone to ``clickbait'' dynamics, where high-entropy, low-value content exploits the user's attention, leading to rapid fatigue, churn, and a breakdown of the latent trust between the user and the digital platform \citep{zhao2023survey, achiam2023gpt}.

The fundamental gap in current advertising research is the failure to model the user's internal world-model. In most industrial systems \citep{touvron2023llama, dubey2024llama}, the user is viewed as a static data source rather than an active agent who navigates the information landscape to minimize their own uncertainty. As the scale of generative content grows \citep{rombach2022high, peebles2023scalable}, the misalignment between the ad system's objective (maximizing surprise to grab attention) and the user's cognitive objective (minimizing surprise to maintain coherence) becomes a critical bottleneck. To address this, we must move beyond the ``passive observation'' of RL and embrace a framework that accounts for the mutual dynamics of perception and action \citep{silver2016mastering, pathak2024foundation}.

In this paper, we propose a paradigm shift from Reward Maximization to \textbf{Free Energy Minimization}. Drawing from the Free Energy Principle (FEP) in theoretical neuroscience \citep{friston2010free}, we introduce the \textbf{Active-Bidder} framework. In this framework, the advertising system and the user are modeled as coupled Active Inference agents. The objective of the advertiser is not merely to ``buy a click,'' but to provide ``Epistemic Value'' that aligns with the user’s internal state. By treating the user's response as a reflection of their surprise relative to a hidden internal model, we can calibrate bidding dynamics to prioritize relevance and long-term utility over short-term diversion. This represents an ontological shift from viewing users as reactive targets to viewing them as active Bayesian navigators \citep{yang2024zero, bronstein2024geometric}.

Mathematically, we replace the standard bid calculation ($b = v \cdot pCTR$) with a Variational Free Energy ($\mathcal{F}$) objective. This objective decomposes into a ``Complexity'' term—measuring the divergence between the system's belief and the user's prior—and an ``Accuracy'' term, representing the relevance of the ad creative to the user's current context \citep{kingma2014adam, goodfellow2014generative}. We define a joint loss function, $\mathcal{L} = \mathcal{F}_{\text{user}} + \lambda \mathcal{F}_{\text{system}}$, which penalizes ``surprisal'' in the interaction loop. Under this regime, an ad that induces high surprise (clickbait) increases the free energy of the system even if it results in a click, thereby reducing its long-term bid value. This calibration ensures that the bidding system naturally filters out intrusive content in favor of ads that provide high information gain relative to the user's hidden states \citep{kovachki2023neural, smith2024scaling}.

To realize this vision at scale, we develop a novel \textbf{Dual-Tower Active Inference Architecture}. The first tower, the Generative Model, predicts the user's expected sensory input (what information they are currently seeking) using state-space modeling techniques \citep{gu2023mamba, vaswani2023evolution}. The second tower, the Recognition Model, encodes the ad creative into the same latent space \citep{radford2021learning}. The bid price is then derived as the Negative Free Energy (the Evidence Lower Bound, or ELBO). To meet the stringent sub-10ms latency requirements of modern RTB, we implement custom CUDA kernels for sparse KL-divergence computation on large-scale embedding tables, leveraging recent advances in hardware-aware attention and training \citep{dao2023flashattention, sun2024learning, reid2024llama}.

We evaluate \textit{Active-Bidder} using the \textit{Social Contract} simulator to model complex user fatigue dynamics and on large-scale real-world datasets including Criteo and Avazu. Our results demonstrate that while traditional RL-based bidders \citep{sutton2018reinforcement} may lead in short-term CTR, \textit{Active-Bidder} significantly outperforms all baselines in long-term User Lifetime Value (LTV) and information-theoretic coherence. This work bridges the gap between biological intelligence and computational auction theory, offering a path toward a more sustainable and aligned digital advertising ecosystem \citep{wu2024ai, wang2024generalist}.

Our contributions are as follows:
\begin{itemize}
    \item We formulate the real-time bidding problem as a \textbf{Joint Free Energy Minimization} task, providing the first rigorous application of Active Inference to computational advertising.
    \item We introduce a \textbf{Dual-Tower Active Inference} architecture that explicitly models user ``surprisal'' and provides a principled mechanism to penalize clickbait and user fatigue.
    \item We develop a \textbf{high-performance CUDA kernel} for sparse variational inference, enabling the deployment of FEP-based models within the millisecond constraints of RTB systems.
    \item We demonstrate through extensive simulation and offline replay \citep{liu2023medtime} that \textit{Active-Bidder} achieves a superior balance between short-term revenue and long-term user retention compared to SOTA RL and LLM-based bidders \citep{zhao2023survey, achiam2023gpt}.
\end{itemize}

The remainder of this paper is organized as follows: Section 2 reviews the foundations of Active Inference and its relation to modern Deep Learning. Section 3 details the mathematical derivation of the Active-Bidder objective. Section 4 describes the architectural implementation and hardware optimizations. Finally, Sections 5 and 6 present our experimental results and a discussion on the ethical alignment of advertising systems.